\section{The Future of Tradition}

There is a reason that I was part of the last ripple. A few years after I had commenced elementary school, the United States Supreme Court (SCOTUS) ruled that any religious expression in public schools was henceforth illegal. I don't believe that many people understand the impact of that decision in its psychological depths. Now, if the legislature had made that decision, it would have been a mere policy decision, without import. That law could itself be overturned and religious education reinstated.

However, the SCOTUS ruling said much more: it said that prayer and religious education in a public school was illegal, moreover, it was a crime. A teacher who started the day with a prayer would suddenly be creating a criminal act, subject to fines and even imprisonment for contempt. A police action would not even be required, but rather a complaint by a student or his parents. Although the teacher could not be punished ex post facto, the meaning is clear. She had been committing a crime and getting away with it. Furthermore, the state of Massachusetts had been a conspirator in a crime for over 300 years.

Furthermore, the notion that the formation of a gentleman required religious instruction was instantly refuted; it was not only not necessary, but even antithetical. The net effect is that the spiritual authority, even in its attenuated form of ``mere Christianity'' as C S Lewis famously put it, was pushed back; henceforth, the political power had pride of place, and not just any political power, but specifically the Federal power. Spiritual authority was restricted to the private sphere.

Obviously, if religious education was implicated in a criminal conspiracy, its entire worldview, including morals and customs, lost their support. Once that conclusion was reached, it did not take much to begin overturning everything that for 300 years had been considered normal and healthy, or, put another way, the established order of things began to be overturned. Now, I am not claiming that people were crypto-Marxists, rather that Marxist was a brilliant philosopher who could see clearly the consequences of rejecting the dominant spiritual authority; the rest just followed along, determined, as it were, by historical forces.

Hence, the movements of so-called conservatives are doomed to failure under such circumstances. No religious or spiritual argument is allowed, and, a fortiori, such arguments are evidence of complicity in a multi-generational crime. For example, the Declaration of Independence mentions rights ``endowed by the Creator''; although the Declaration is part of the organic law of the USA, that argument would not be allowed in any court. One of those rights is the ``right to life''. A derivative right, therefore, is the right to self-defense. And the right to bear arms is clearly necessary for self-defense. That is how the argument should go, rather than on the merely legal basis of the 2\textsuperscript{nd} Amendment.

In certain circles, the right to bear arms is under attack and is considered déclassé. If so, then also is the right to self-defense and then the right to life. Is this far fetched? Only if you do not know history. Herodotus wrote about King Cyrus' policy toward the Lydians:

\begin{quotex}
Send a message that they are forbidden to own weapons of war, that they are to wear tunics under their coats and slippers on their feet, they are to take up the cithara and the harp, and that they are to raise their sons to be retailers. Before long you will see them become women instead of men, and so there will be no danger of them rising up against you.
\end{quotex}


The Lydians today are those who still do not see the ways of their ancestors as criminal, ignorant, or immoral, despite the weight of the court's rulings. There is now no danger of them rising up against anyone.

\paragraph{Thelema, the Pope, and Chaos}


I'm sure I must be misunderstanding something in this recent interview of Pope Francis\footnote{\url{https://rorate-caeli.blogspot.com/2013/10/the-pope-to-repubblica-how-church-must.html}}, it seems to me that Pope Francis is saying, ``Love and do what you will,''\footnote{That quote is actually from St Augustine, one of the Pope's favorite saints. Good choice.} i.e., belief and dogma are unnecessary. In the interview, there is a lot of talk of social justice, but prayer or meditation were never mentioned. He considers everyone a child of God, much like Oprah Winfrey, but I do not believe that is dogmatic or scriptural. Only the baptized are children of God, and some are even called the offspring of Satan by Jesus.

What is more interesting to us is the belief of the interviewer Eugenio Scalfari. He believes in ``Being'', particularly that Being is chaotic energy. Like all pagans, Being is Chaos and form is the accidental result of these energies combining in certain ways.

Now ``Being'' is the God that the Pope mentions several times. The Pope also believes in Jesus Christ, the ``incarnation of God'', or more specifically, the incarnation of the Logos. Despite the frequent call to engage with the modern world, this was an opportunity missed. The fundamental decision is this: is Being united to Chaos or to Logos? That, it seems to me, is a question that Scalfari could have understood. These are the dogmas of the tradition of the future, at least version 1.0:

\begin{enumerate}


\item God is the ground of Being

\item Being is united to and co-eternal with Logos

\item Logos is the foundational source of creative action, and Chaos is not real in itself, but is rather the privation of Logos

\item Material forms are not accidental, since they are condensed energy

\item Energy is not chaotic, since it is condensed Will

\item Will is not the same as desire, since it arises from Intelligence

\item Intelligence is guided by Love, the force of attraction, which leads us toward transcendence
\end{enumerate}



\flrightit{Posted on 2013-10-02 by Cologero}

\begin{center}* * *\end{center}

\begin{footnotesize}
\begin{sffamily}
\texttt{Seraphim on 2013-10-06 at 22:31 said:}

How do you get ``belief and dogma are unnecessary'' from St. Augustine's ``Love God, and do what you will''?

We cannot separate love, and contemplatio/Intellect. They do not exist independently of each other.

\hfill

\texttt{Cologero on 2013-10-06 at 22:44 said:}

Forgive my unclear prose, Seraphim, the point was supposed to be that many, if not most, will come, or actually have come, to that conclusion.


\hfill

\end{sffamily}
\end{footnotesize}

